\chapter{农业、建筑、巡检与危险环境}
\label{chap:field-dangerous-environment-case}

\section{现场任务由环境与通信共同定义}

农业、建筑、矿山、能源巡检和灾害响应的共同点不是“户外”，而是地形、光照、粉尘、水、温度、通信和维护条件难以固定。实验室成功的导航器若依赖连续网络、洁净镜头或人工复位，在地下现场可能立刻失效。本章以 DARPA Subterranean Challenge 的 CERBERUS 系统为完整案例，不再逐项罗列行业。

DARPA 将任务定义为在未知地下环境中快速建图、导航、搜索并定位目标\cite{darpa2021subt}。竞赛同时引入隧道、城市地下和洞穴环境，把通信退化、垂直结构、黑暗、粉尘和无 GPS 变成系统约束。它仍是规则化竞赛，不代表矿山或消防的全部法规、耐久和人员流程。

\section{CERBERUS：异构系统而非单一最强机器人}

CERBERUS 组合 ANYmal 腿式机器人、常规与耐碰撞飞行器、多模态视觉/LiDAR/惯性定位、探索规划、目标检测、集中地图优化和通信中继，在 2021 决赛取得胜利\cite{tranzatto2022cerberus}。Tunnel/Urban Circuits 的系统论文记录了现场部署与限制\cite{tranzatto2022tunnel}。

\begin{figure}[H]
\centering
\begin{tabular}{c@{\ $\leftrightarrow$\ }c@{\ $\leftrightarrow$\ }c}
\fbox{腿式/飞行探索机器人} & \fbox{中继、投放通信节点} & \fbox{基站、操作员与地图融合}
\end{tabular}

\vspace{0.6em}
\begin{tabular}{c@{\quad}c@{\quad}c}
\fbox{局部自治与避障} & \fbox{失联继续/返航} & \fbox{任务分配与人工监督}
\end{tabular}
\caption{地下异构系统的通信与控制结构。自主性由机器人和基站共同组织。}
\label{fig:cerberus-system-case}
\end{figure}

腿式平台越障能力强但续航有限；飞行器进入狭窄/高处但通信和能量更紧；地面中继与投放通信节点扩展网络；集中优化改善多机器人地图一致性。异构性的代价是备件、软件、标定、操作和故障模式增加，因此统一的任务/地图接口比统一本体更重要。

任务效用可抽象为
\begin{equation}
J=w_c C+w_d N_{\mathrm{detected}}-w_e E-w_h H-w_r R,
\label{eq:field-mission-utility}
\end{equation}
其中 $C$ 是覆盖或信息增益，$N_{\mathrm{detected}}$ 是正确定位目标数，$E$ 是能耗，$H$ 是人工负担，$R$ 是机器人丢失或安全风险。竞赛分数主要显式奖励目标和时间，真实运营还需提高风险、耐久和人员安全权重。

\section{完全自主、遥操作与共享自主}

\begin{table}[H]
\begingroup\hbadness=10000
\centering
\caption{现场自治组织方式的适用条件}
\label{tab:field-autonomy-modes}
\begin{tabularx}{\textwidth}{p{0.18\textwidth}YYYp{0.20\textwidth}}
\toprule
模式 & 优势 & 脆弱点 & 合适任务 & 必报指标 \\
\midrule
完全自主 & 可在失联区持续工作 & 未知场景判断与恢复困难 & 高重复、边界清楚、低后果 & 无人时长、失败自恢复 \\
连续遥操作 & 人直接理解复杂语境 & 带宽、时延、操作者负荷 & 短程高价值精细操作 & 链路可用率、人员分钟 \\
共享自主 & 机器执行局部闭环，人管目标/异常 & 接管界面与责任切换复杂 & 地下探索、巡检、应急 & 介入原因、接管延迟、恢复率 \\
多机器人监督 & 一人覆盖更大区域 & 告警洪泛与地图冲突 & 搜索、测绘、巡检 & 人机比、覆盖、冲突和丢机 \\
\bottomrule
\end{tabularx}
\endgroup
\end{table}

CERBERUS 的“自主”并非没有人：操作员在基站分配任务、监控地图和机器人健康，并在需要时干预；机器人在局部执行探索、规划和避障。现场可扩展性的关键指标是每名操作员能可靠监督多少机器人、每小时需要多少干预，而不是是否存在人工。

\section{现场约束如何反向塑造系统}

\begin{table}[H]
\begingroup\hbadness=10000
\centering
\caption{环境约束到技术机制的映射}
\label{tab:field-constraint-mechanism}
\begin{tabularx}{\textwidth}{p{0.20\textwidth}YYp{0.23\textwidth}}
\toprule
现场约束 & 设计机制 & 失败检测 & 运维证据 \\
\midrule
无 GPS、退化视觉 & LiDAR/视觉/IMU 融合与地图优化 & 协方差、回环冲突、漂移 & 轨迹真值片段、重定位率 \\
间歇通信 & 本地自治、中继、存储转发 & 心跳、队列、链路质量 & 覆盖图、失联时长、数据丢失 \\
崎岖/狭窄地形 & 腿式与飞行异构平台 & 足端滑移、碰撞、卡滞 & 通过率、能耗、损伤 \\
有限续航 & 信息增益规划、返航阈值 & 电池、温度、剩余路径 & 有效任务时长、换电工时 \\
高维护成本 & 模块化载荷、远程诊断、备件 & 健康监测与故障码 & MTBF、MTTR、现场工具 \\
\bottomrule
\end{tabularx}
\endgroup
\end{table}

农业还会增加季节与生物变化，建筑会增加动态工序与临时结构，工业巡检会增加防爆和停产窗口。它们共享现场方法，但不能从 SubT 成绩直接外推。每个新行业都要重新定义危险、通信、耐久、清洁、责任和单位任务收益。

\section{从竞赛证据到长期运营还差什么}

SubT 提供公开规则、未知现场、时间压力和跨团队比较，证据强于剪辑 demo；但赛事持续时间短，团队可集中专家，机器人损失和维护成本与商业现场不同。运营需要月级/年级暴露量、天气和季节覆盖、备件、远程支持、培训、事故与近失事件、数据回传和合规记录。

商业模式也受自治组织反向牵引。若通信和异常要求专家长期在线，系统更像远程服务而非硬件销售；若一名操作员能监督多机且维护标准化，RaaS 才可能摊薄成本；若任务低频但危险，价值可能来自减少人员暴露而非提高吞吐。

\section{知识小结与案例判断}

CERBERUS 的现场能力来自异构本体、多模态定位、探索规划、通信工程、地图融合和人工监督共同作用。竞赛证明了受规则约束的高强度系统集成，不证明长期商业运营。案例判断是：危险环境最现实的路线通常是有界局部自治加共享监督；评价必须同时报告覆盖、目标、失联、能耗、人员负担、丢机和维护，而不能只给“自主完成”标签。
